% Vietnamese translation for OI Public Library
% Source-PDF: 数据结构 DATA STRUCTURE/后缀自动机 - 陈立杰.pdf
\documentclass[11pt]{article}
\usepackage{oipl-vn}

\title{Máy tự động hậu tố}
\author{Trần Lập Kiệt, Trường Ngoại ngữ Hàng Châu}
\date{}

\begin{document}
\maketitle

\origtitle{Suffix Automaton}

\section{Mở đầu}

Tài liệu này giới thiệu \term{máy tự động hậu tố} (suffix automaton, viết tắt là
SAM) và một số ứng dụng kinh điển trong xử lý xâu.  Mục tiêu là có một cấu trúc
tuyến tính theo độ dài xâu nhưng vẫn trả lời được nhiều câu hỏi về xâu con,
vị trí xuất hiện và thứ tự từ điển.

Trước hết xét bài SPOJ 1812, \term{Longest Common Substring II}: cho
$N \le 10$ xâu, mỗi xâu có độ dài không quá $100000$, hãy tìm xâu con liên tiếp
chung dài nhất của tất cả các xâu.  Một cách đơn giản là nhị phân độ dài rồi
dùng băm, đạt $O(L\log L)$ với $L$ là tổng độ dài.  Cách này quen thuộc nhưng
trên môi trường chậm như SPOJ vẫn dễ bị quá thời gian.  SAM đưa ra một hướng
khác, tuyến tính hơn và phù hợp với các bài trực tuyến.

\section{Các công cụ xử lý xâu trong OI}

Trong OI, các công cụ xử lý xâu thường gặp gồm:
\begin{itemize}
  \item mảng hậu tố (suffix array);
  \item cây hậu tố (suffix tree);
  \item máy tự động Aho--Corasick;
  \item băm xâu;
  \item máy tự động hậu tố.
\end{itemize}

SAM có thể xem là một phiên bản nén của cây trie các hậu tố, nhưng số trạng thái
và số cạnh đều chỉ tuyến tính.  So với cây hậu tố, thuật toán xây dựng của SAM
ngắn và dễ cài đặt hơn, dù các bất biến ban đầu có thể khó nhìn.

\section{Máy tự động hữu hạn}

Một máy tự động hữu hạn dùng để nhận dạng xâu.  Với máy tự động $A$, nếu $A$
nhận xâu $S$ thì viết $A(S)=\mathrm{True}$, ngược lại là $\mathrm{False}$.

Một máy tự động gồm năm phần:
\begin{itemize}
  \item $\Sigma$: bảng chữ cái;
  \item $\mathrm{state}$: tập trạng thái;
  \item $\mathrm{init}$: trạng thái ban đầu;
  \item $\mathrm{end}$: tập trạng thái kết thúc;
  \item $\mathrm{trans}$: hàm chuyển trạng thái.
\end{itemize}

Ký hiệu $\mathrm{trans}(s,c)$ là trạng thái đi đến khi đang ở trạng thái $s$ và
đọc ký tự $c$.  Nếu chuyển này không tồn tại, ta xem kết quả là
$\mathrm{null}$; trạng thái $\mathrm{null}$ chỉ chuyển về chính nó.  Mở rộng
cho xâu $t$, $\mathrm{trans}(s,t)$ là trạng thái thu được sau khi lần lượt đọc
từng ký tự của $t$:
\begin{verbatim}
cur = s
for i = 0 .. length(t)-1:
    cur = trans(cur, t[i])
\end{verbatim}

\section{Định nghĩa SAM}

Cho xâu mẹ $S$.  Máy tự động hậu tố của $S$ là một máy tự động nhận đúng tất cả
các hậu tố của $S$:
\[
  \mathrm{SAM}(x)=\mathrm{True}
  \quad\Longleftrightarrow\quad
  x \text{ là hậu tố của } S.
\]

Sau khi xây dựng xong, cùng cấu trúc này cũng nhận biết được tất cả xâu con của
$S$.  Ký hiệu
\[
  \mathrm{ST}(t)=\mathrm{trans}(\mathrm{init},t)
\]
là trạng thái đạt được khi đọc $t$ từ gốc.  Khi $\mathrm{ST}(t)\ne
\mathrm{null}$, $t$ là một xâu con của $S$.

Trong ký hiệu của slide gốc, \(\mathrm{Suf}\) là tập hậu tố của \(S\),
\(\mathrm{Fac}\) là tập mọi xâu con liên tiếp của \(S\), \(\mathrm{Suffix}(i)\)
là hậu tố bắt đầu tại vị trí \(i\), và \(S[l,r]\) là xâu con trên đoạn
\([l,r]\) với chỉ số bắt đầu từ \(0\). Nếu \(t\notin\mathrm{Fac}\) thì
\(\mathrm{ST}(t)=\mathrm{null}\), vì nối thêm bất kỳ phần sau nào cũng không thể
biến \(t\) thành hậu tố của \(S\). Ngược lại, nếu \(t\in\mathrm{Fac}\) thì
\(\mathrm{ST}(t)\ne\mathrm{null}\), vì có thể nối thêm phần đuôi sau một lần
xuất hiện của \(t\) để tạo thành một hậu tố của \(S\).

Cách thô sơ nhất là đưa mọi hậu tố của $S$ vào một trie.  Khi đó gốc là trạng
thái ban đầu, cạnh trie là hàm chuyển, còn các lá là trạng thái kết thúc.  Tuy
nhiên trie này có thể có kích thước $O(n^2)$.  SAM tối tiểu hóa các trạng thái
có cùng hành vi, nhờ đó chỉ còn $O(n)$ trạng thái.

\section{Tập vị trí kết thúc và trạng thái}

Với một xâu con $t$ của $S$, đặt
\[
  \mathrm{Right}(t)=\{i \mid S[i-|t|+1..i]=t\},
\]
tức là tập mọi vị trí kết thúc của những lần xuất hiện của $t$ trong $S$.

Trong SAM tối tiểu, một trạng thái đại diện cho một lớp các xâu con có cùng tập
$\mathrm{Right}$.  Nếu trạng thái là $s$, ta viết tập này là
$\mathrm{Right}(s)$.  Các độ dài của xâu con thuộc cùng trạng thái tạo thành
một đoạn liên tiếp:
\[
  [\mathrm{Min}(s),\mathrm{Max}(s)].
\]
Một cặp gồm tập $\mathrm{Right}$ và độ dài xác định duy nhất một xâu con.

Với mỗi trạng thái $s$ khác gốc, $\mathrm{parent}(s)$ là trạng thái $x$ sao cho
$\mathrm{Right}(s)$ là tập con thực sự của $\mathrm{Right}(x)$, và
$|\mathrm{Right}(x)|$ nhỏ nhất có thể.  Các liên kết $\mathrm{parent}$ tạo thành
một cây, thường gọi là \term{cây parent} hay \term{suffix-link tree}.  Ta có bất
biến quan trọng:
\[
  \mathrm{Max}(\mathrm{parent}(s))=\mathrm{Min}(s)-1.
\]

\section{Tính chất về xâu con}

Mọi xâu con của $S$ đều nằm trong một trạng thái nào đó của SAM.  Do đó:
\[
  t \text{ là xâu con của } S
  \quad\Longleftrightarrow\quad
  \mathrm{ST}(t)\ne \mathrm{null}.
\]

Nếu đã đi đến trạng thái chứa $t$, số lần xuất hiện của $t$ chính là
$|\mathrm{Right}(t)|$, tức kích thước của tập $\mathrm{Right}$ tương ứng.

Không nên lưu trực tiếp toàn bộ tập $\mathrm{Right}$ ở mỗi trạng thái vì tốn bộ
nhớ.  Ta dùng cây parent: tập $\mathrm{Right}$ của một trạng thái là hợp các tập
$\mathrm{Right}$ của những lá trong cây con của nó.  Nếu đánh số theo thứ tự DFS
của cây parent, toàn bộ một cây con là một đoạn liên tiếp; nhờ đó có thể dùng
các cấu trúc dữ liệu trên đoạn để tìm nhanh mọi vị trí xuất hiện của một xâu
con.

\section{Xây dựng tuyến tính}

Thuật toán xây dựng SAM là trực tuyến: đọc xâu từ trái sang phải và thêm từng ký
tự.  Dù ý tưởng của nó không hiển nhiên, phần cài đặt ngắn hơn nhiều so với cây
hậu tố.

Giả sử SAM hiện tại biểu diễn xâu $S$ có độ dài $L$, trạng thái
$\mathrm{last}$ đại diện cho toàn bộ xâu $S$.  Khi thêm ký tự $c$, ta tạo trạng
thái mới $np$ với
\[
  \mathrm{Max}(np)=L+1.
\]
Trạng thái này đại diện cho các hậu tố mới kết thúc ở vị trí $L+1$.

Từ $p=\mathrm{last}$, đi ngược theo các liên kết parent.  Với mọi trạng thái
$p$ chưa có cạnh $c$, thêm cạnh $p \xrightarrow{c} np$.  Khi quá trình dừng có
ba trường hợp.

\begin{enumerate}
  \item Không còn trạng thái nào, tức đã đi qua gốc.  Khi đó
  $\mathrm{parent}(np)=\mathrm{init}$.

  \item Tồn tại cạnh $p \xrightarrow{c} q$ và
  $\mathrm{Max}(p)+1=\mathrm{Max}(q)$.  Khi đó $q$ đã là trạng thái đúng để nối
  tiếp, đặt $\mathrm{parent}(np)=q$.

  \item Tồn tại cạnh $p \xrightarrow{c} q$ nhưng
  $\mathrm{Max}(p)+1<\mathrm{Max}(q)$.  Khi đó cần tách $q$.  Tạo trạng thái
  sao chép $nq$, sao chép toàn bộ chuyển trạng thái của $q$, đặt
  $\mathrm{Max}(nq)=\mathrm{Max}(p)+1$ và
  $\mathrm{parent}(nq)=\mathrm{parent}(q)$ cũ.  Sau đó đặt
  $\mathrm{parent}(q)=\mathrm{parent}(np)=nq$, rồi đổi mọi cạnh
  $v \xrightarrow{c} q$ trên đường parent liên quan thành
  $v \xrightarrow{c} nq$.
\end{enumerate}

Cuối cùng đặt $\mathrm{last}=np$.  Trong bước tách, các chuyển của $nq$ giống
$q$ vì khi đang xử lý ký tự mới ở vị trí $L+1$, phần sau trạng thái này không
phụ thuộc riêng vào vị trí mới.

\section{Mã giả}

\begin{verbatim}
extend(c):
    np = new_state()
    maxlen[np] = maxlen[last] + 1
    p = last

    while p exists and next[p][c] is absent:
        next[p][c] = np
        p = parent[p]

    if p does not exist:
        parent[np] = root
    else:
        q = next[p][c]
        if maxlen[p] + 1 == maxlen[q]:
            parent[np] = q
        else:
            nq = clone_state(q)
            maxlen[nq] = maxlen[p] + 1
            parent[nq] = parent[q]
            parent[q] = nq
            parent[np] = nq
            while p exists and next[p][c] == q:
                next[p][c] = nq
                p = parent[p]

    last = np
\end{verbatim}

Mỗi lần thêm ký tự chỉ tạo nhiều nhất hai trạng thái.  Vì vậy số trạng thái của
SAM không vượt quá $2n-1$ với $n>0$.  Số cạnh cũng tuyến tính; nhìn trực quan,
các cạnh được thêm và đổi hướng dọc theo các suffix link, tổng số lần thao tác
trên toàn bộ quá trình là tuyến tính.

\section{I. Xâu quay vòng nhỏ nhất}

Cho xâu $S$.  Mỗi lần được chuyển ký tự đầu tiên xuống cuối xâu.  Hãy tìm xâu
có thứ tự từ điển nhỏ nhất có thể thu được.  Ví dụ với \texttt{BBAAB}, kết quả
là \texttt{AABBB}.

Xét xâu $SS$.  Bài toán trở thành tìm xâu con độ dài $|S|$ nhỏ nhất theo thứ tự
từ điển của $SS$.  Xây SAM cho $SS$; từ trạng thái gốc, mỗi lần đi theo cạnh có
nhãn nhỏ nhất, đi đúng $|S|$ bước thì thu được đáp án.  Lý do là mọi phép quay
của $S$ đều là một xâu con độ dài $|S|$ trong $SS$, và đi tham lam theo ký tự
nhỏ nhất sinh ra xâu nhỏ nhất trong ngôn ngữ được SAM biểu diễn.

\section{II. SPOJ NSUBSTR}

Cho xâu $S$.  Gọi $F(x)$ là số lần xuất hiện lớn nhất trong tất cả xâu con độ
dài $x$ của $S$.  Cần tính $F(1),F(2),\ldots,F(|S|)$, với
$|S|\le 250000$.

Xây SAM của $S$.  Với mỗi trạng thái $s$, các xâu mà nó đại diện có độ dài nằm
trong đoạn
\[
  [\mathrm{Min}(s),\mathrm{Max}(s)],
\]
và cùng có số lần xuất hiện $|\mathrm{Right}(s)|$.  Vì thế dùng
$|\mathrm{Right}(s)|$ để cập nhật $F(\mathrm{Max}(s))$.  Sau đó duyệt $i$ giảm
dần và gán
\[
  F(i-1)=\max(F(i-1),F(i)).
\]
Việc lan truyền này đúng vì nếu một xâu dài xuất hiện nhiều lần thì tiền tố
ngắn hơn của nó cũng xuất hiện ít nhất bấy nhiêu lần.

\section{III. BZOJ2555 SubString}

Bài toán yêu cầu duy trì một xâu, hỗ trợ hai thao tác trực tuyến:
\begin{itemize}
  \item thêm một ký tự vào cuối xâu;
  \item hỏi một xâu xuất hiện bao nhiêu lần trong xâu hiện tại.
\end{itemize}

Ta xây SAM động bằng thủ tục thêm ký tự ở trên.  Điều thực sự ảnh hưởng đến đáp
án là kích thước tập $\mathrm{Right}$ của trạng thái truy vấn.  Trong mỗi bước
xây dựng, liên kết parent chỉ thay đổi hằng số lần; trạng thái mới $np$ tương
ứng với một lần xuất hiện mới, nên mọi tổ tiên của $np$ trên cây parent cần tăng
kích thước $\mathrm{Right}$ thêm $1$.

Có hai hướng duy trì:
\begin{itemize}
  \item dùng cây động để duy trì cây parent;
  \item dùng cây cân bằng để duy trì thứ tự DFS của cây parent.
\end{itemize}
Các kỹ thuật này nằm ngoài trọng tâm của bài giảng, nên tài liệu gốc chỉ nêu ý
tưởng.

\section{IV. SPOJ SUBLEX}

Cho một xâu $S$ độ dài không quá $90000$.  Mỗi truy vấn hỏi xâu thứ $K$ theo
thứ tự từ điển trong tập các xâu con phân biệt của $S$; số truy vấn không quá
$500$.

Xây SAM của $S$.  Với mỗi trạng thái, tiền xử lý số lượng xâu con phân biệt có
thể đi tới từ trạng thái đó.  Sau đó trả lời truy vấn bằng cách đi từ gốc theo
thứ tự nhãn cạnh tăng dần: nếu toàn bộ nhánh hiện tại chứa ít hơn $K$ xâu thì bỏ
qua nhánh đó và trừ vào $K$; ngược lại chọn cạnh đó, in thêm ký tự cạnh, rồi
tiếp tục.  Đây là cách DFS theo thứ tự từ điển trên DAG của SAM.

\section{V. SPOJ LCS}

Cho hai xâu $A$ và $B$, mỗi xâu có độ dài nhỏ hơn $100000$, hãy tìm xâu con
liên tiếp chung dài nhất.  Xây SAM cho $A$.

Duyệt các ký tự của $B$ bằng SAM của $A$.  Duy trì trạng thái hiện tại $s$ và
độ dài khớp hiện tại $\ell$.
\begin{itemize}
  \item Nếu có cạnh nhãn $c$ từ $s$, đi theo cạnh đó và tăng $\ell$.
  \item Nếu không có, đi ngược theo parent cho đến khi gặp trạng thái $a$ có
  cạnh nhãn $c$.  Khi đó chuyển đến $\mathrm{trans}(a,c)$ và đặt
  $\ell=\mathrm{Max}(a)+1$.
  \item Nếu không có tổ tiên nào có cạnh nhãn $c$, quay về gốc và đặt
  $\ell=0$.
\end{itemize}

Trong quá trình duyệt, cập nhật độ dài khớp lớn nhất đạt được ở từng trạng
thái.  Vì kết quả của một trạng thái cũng có thể đóng góp cho parent của nó,
cuối cùng duyệt các trạng thái theo $\mathrm{Max}$ giảm dần để đẩy kết quả lên
cây parent.  Đáp án là giá trị lớn nhất thu được.

\section{VI. SPOJ LCS2}

Bài LCS2 mở rộng bài trước cho nhiều xâu.  Xây SAM của xâu đầu tiên.  Với mỗi
xâu còn lại, chạy quy trình khớp như trong SPOJ LCS để biết tại mỗi trạng thái
độ dài lớn nhất mà xâu này có thể khớp.  Sau khi xử lý một xâu, lan truyền kết
quả lên cây parent, rồi lấy minimum theo từng trạng thái qua tất cả các xâu.
Giá trị lớn nhất còn lại chính là độ dài xâu con chung dài nhất của toàn bộ tập
xâu.

\section{Một vài cấu trúc liên quan}

Ngoài suffix automaton còn có nhiều biến thể và họ hàng:
\begin{itemize}
  \item factor automaton;
  \item suffix oracle;
  \item factor oracle;
  \item suffix cactus.
\end{itemize}

Factor oracle là một máy tự động có thể nhận tất cả xâu con của xâu đã cho,
nhưng cũng có thể nhận thêm một số xâu không thuộc tập đó.  Vì vậy chữ
``oracle'' ở đây đúng nghĩa là một lời dự đoán: nhanh và đơn giản, nhưng không
nhất thiết chính xác tuyệt đối.  Tài liệu gốc chỉ nhắc tới nó như phần mở rộng,
không đi sâu vào ứng dụng trong OI.

\section{Ghi chú về bản dịch}

Một số trang của PDF là hình minh họa hoặc mã nhúng dưới dạng đối tượng đồ
họa, nên không có lớp văn bản để trích xuất trực tiếp.  Bản dịch vẫn giữ trình
tự bài giảng và phục dựng các ý thuật toán chính: định nghĩa
$\mathrm{Right}$, cây parent, thuật toán thêm ký tự, và các ứng dụng nêu trong
slide.

\end{document}
